%%%%%%%%%%%%%%%%%%%%%%%%%%%%%%%%%%%%%%%%%%%%%%%%%%%%%%%%%%%%
% 5.10 FOURIER ANALYSIS
% The Composition of Advanced Mathematics
%%%%%%%%%%%%%%%%%%%%%%%%%%%%%%%%%%%%%%%%%%%%%%%%%%%%%%%%%%%%

\section{Fourier Analysis}
\label{sec:fourier-analysis}

\prerequisite{
Trigonometric functions, complex exponentials, sequences and series,
integration, improper integrals, complex analysis, measure theory,
Lebesgue integration, \(L^p\) spaces, inner products, Hilbert spaces,
and basic functional analysis.
}

\learningobjectives{
By the end of this section, the reader should be able to:
\begin{itemize}
    \item explain the central idea of frequency decomposition;
    \item compute real and complex Fourier coefficients;
    \item construct Fourier series of periodic functions;
    \item identify even and odd symmetry in Fourier expansions;
    \item work with sine and cosine series;
    \item interpret Fourier series as orthogonal expansions in \(L^2\);
    \item apply Bessel's inequality and Parseval's identity;
    \item distinguish pointwise, uniform, and \(L^2\) convergence;
    \item understand the Dirichlet kernel and partial-sum operators;
    \item recognize the Gibbs phenomenon;
    \item define and use convolution;
    \item understand approximate identities and Fej\'er summation;
    \item define the Fourier transform on \(\R\);
    \item use translation, modulation, scaling, and differentiation rules;
    \item apply the convolution theorem;
    \item understand the Riemann--Lebesgue Lemma;
    \item use Plancherel's Theorem and Fourier inversion;
    \item work with the Fourier transform on \(L^2\);
    \item understand the uncertainty principle;
    \item recognize the role of Fourier methods in differential equations,
          signal processing, probability, and harmonic analysis.
\end{itemize}
}

%%%%%%%%%%%%%%%%%%%%%%%%%%%%%%%%%%%%%%%%%%%%%%%%%%%%%%%%%%%%
\subsection{The Central Idea of Fourier Analysis}
\label{subsec:central-idea-fourier}
%%%%%%%%%%%%%%%%%%%%%%%%%%%%%%%%%%%%%%%%%%%%%%%%%%%%%%%%%%%%

Fourier analysis studies functions by decomposing them into elementary
oscillatory components.

The fundamental building blocks are

\[
\boxed{
\cos(\omega x),
\qquad
\sin(\omega x),
}
\]

or, equivalently,

\[
\boxed{
e^{i\omega x}.
}
\]

The parameter \(\omega\), lowercase Greek \emph{omega}, represents angular
frequency.

The central principle is

\[
\boxed{
\text{complicated function}
\longleftrightarrow
\text{collection of frequencies}.
}
\]

This allows a problem in the original variable to be transformed into a
problem involving frequency.

%%%%%%%%%%%%%%%%%%%%%%%%%%%%%%%%%%%%%%%%%%%%%%%%%%%%%%%%%%%%
\subsection{Periodic Functions}
\label{subsec:periodic-functions-fourier}
%%%%%%%%%%%%%%%%%%%%%%%%%%%%%%%%%%%%%%%%%%%%%%%%%%%%%%%%%%%%

A function \(f\) is periodic with period \(T>0\) if

\[
\boxed{
f(x+T)=f(x)
}
\]

for every \(x\) in its domain.

The fundamental angular frequency is

\[
\boxed{
\omega_0=\frac{2\pi}{T}.
}
\]

The corresponding harmonics are

\[
\boxed{
\omega_n=n\omega_0
=
\frac{2\pi n}{T}.
}
\]

Fourier series decompose periodic functions into these discrete
frequencies.

%%%%%%%%%%%%%%%%%%%%%%%%%%%%%%%%%%%%%%%%%%%%%%%%%%%%%%%%%%%%
\subsection{Orthogonality of Trigonometric Functions}
\label{subsec:trigonometric-orthogonality}
%%%%%%%%%%%%%%%%%%%%%%%%%%%%%%%%%%%%%%%%%%%%%%%%%%%%%%%%%%%%

On the interval

\[
[-\pi,\pi],
\]

the functions

\[
1,
\qquad
\cos(nx),
\qquad
\sin(nx)
\]

satisfy orthogonality relations.

For positive integers \(m,n\),

\[
\boxed{
\int_{-\pi}^{\pi}
\cos(mx)\cos(nx)\,dx
=
\begin{cases}
0, & m\neq n,\\
\pi, & m=n,
\end{cases}
}
\]

and

\[
\boxed{
\int_{-\pi}^{\pi}
\sin(mx)\sin(nx)\,dx
=
\begin{cases}
0, & m\neq n,\\
\pi, & m=n.
\end{cases}
}
\]

Also,

\[
\boxed{
\int_{-\pi}^{\pi}
\sin(mx)\cos(nx)\,dx
=
0.
}
\]

These relations are the infinite-dimensional analogue of orthogonal
vectors in Euclidean space.

%%%%%%%%%%%%%%%%%%%%%%%%%%%%%%%%%%%%%%%%%%%%%%%%%%%%%%%%%%%%
\subsection{Real Fourier Series}
\label{subsec:real-fourier-series}
%%%%%%%%%%%%%%%%%%%%%%%%%%%%%%%%%%%%%%%%%%%%%%%%%%%%%%%%%%%%

A \(2\pi\)-periodic function is formally represented by

\[
\boxed{
f(x)
\sim
\frac{a_0}{2}
+
\sum_{n=1}^{\infty}
\left(
a_n\cos(nx)
+
b_n\sin(nx)
\right).
}
\]

The symbol

\[
\sim
\]

indicates a Fourier expansion; convergence must be analyzed separately.

The Fourier coefficients are

\[
\boxed{
a_0
=
\frac{1}{\pi}
\int_{-\pi}^{\pi}f(x)\,dx,
}
\]

\[
\boxed{
a_n
=
\frac{1}{\pi}
\int_{-\pi}^{\pi}
f(x)\cos(nx)\,dx,
}
\]

and

\[
\boxed{
b_n
=
\frac{1}{\pi}
\int_{-\pi}^{\pi}
f(x)\sin(nx)\,dx.
}
\]

%%%%%%%%%%%%%%%%%%%%%%%%%%%%%%%%%%%%%%%%%%%%%%%%%%%%%%%%%%%%
\subsection{Fourier Series on a General Interval}
\label{subsec:fourier-general-interval}
%%%%%%%%%%%%%%%%%%%%%%%%%%%%%%%%%%%%%%%%%%%%%%%%%%%%%%%%%%%%

For a function of period \(2L\), one writes

\[
\boxed{
f(x)
\sim
\frac{a_0}{2}
+
\sum_{n=1}^{\infty}
\left[
a_n
\cos\left(\frac{n\pi x}{L}\right)
+
b_n
\sin\left(\frac{n\pi x}{L}\right)
\right].
}
\]

The coefficients are

\[
\boxed{
a_n
=
\frac1L
\int_{-L}^{L}
f(x)
\cos\left(\frac{n\pi x}{L}\right)
\,dx,
}
\]

and

\[
\boxed{
b_n
=
\frac1L
\int_{-L}^{L}
f(x)
\sin\left(\frac{n\pi x}{L}\right)
\,dx.
}
\]

For \(n=0\),

\[
\boxed{
a_0
=
\frac1L
\int_{-L}^{L}f(x)\,dx.
}
\]

%%%%%%%%%%%%%%%%%%%%%%%%%%%%%%%%%%%%%%%%%%%%%%%%%%%%%%%%%%%%
\subsection{Even Functions}
\label{subsec:even-fourier-series}
%%%%%%%%%%%%%%%%%%%%%%%%%%%%%%%%%%%%%%%%%%%%%%%%%%%%%%%%%%%%

If

\[
f(-x)=f(x),
\]

then \(f\) is even.

Because

\[
\sin(nx)
\]

is odd,

\[
f(x)\sin(nx)
\]

is odd.

Therefore

\[
\boxed{
b_n=0.
}
\]

An even function has a cosine Fourier series:

\[
\boxed{
f(x)
\sim
\frac{a_0}{2}
+
\sum_{n=1}^{\infty}
a_n\cos(nx).
}
\]

Moreover,

\[
\boxed{
a_n
=
\frac{2}{\pi}
\int_0^\pi
f(x)\cos(nx)\,dx.
}
\]

%%%%%%%%%%%%%%%%%%%%%%%%%%%%%%%%%%%%%%%%%%%%%%%%%%%%%%%%%%%%
\subsection{Odd Functions}
\label{subsec:odd-fourier-series}
%%%%%%%%%%%%%%%%%%%%%%%%%%%%%%%%%%%%%%%%%%%%%%%%%%%%%%%%%%%%

If

\[
f(-x)=-f(x),
\]

then \(f\) is odd.

Since

\[
f(x)\cos(nx)
\]

is odd,

\[
\boxed{
a_n=0.
}
\]

An odd function has a sine Fourier series:

\[
\boxed{
f(x)
\sim
\sum_{n=1}^{\infty}
b_n\sin(nx).
}
\]

The coefficients simplify to

\[
\boxed{
b_n
=
\frac{2}{\pi}
\int_0^\pi
f(x)\sin(nx)\,dx.
}
\]

%%%%%%%%%%%%%%%%%%%%%%%%%%%%%%%%%%%%%%%%%%%%%%%%%%%%%%%%%%%%
\subsection{Worked Example: Fourier Series of \(f(x)=x\)}
\label{subsec:worked-fourier-x}
%%%%%%%%%%%%%%%%%%%%%%%%%%%%%%%%%%%%%%%%%%%%%%%%%%%%%%%%%%%%

\begin{workedexamplebox}{Fourier Series of \(f(x)=x\)}

Consider

\[
f(x)=x,
\qquad
-\pi<x<\pi,
\]

extended periodically.

Since \(f\) is odd,

\[
a_0=a_n=0.
\]

Thus

\[
b_n
=
\frac{2}{\pi}
\int_0^\pi
x\sin(nx)\,dx.
\]

Integration by parts gives

\[
\int_0^\pi
x\sin(nx)\,dx
=
-\frac{\pi(-1)^n}{n}.
\]

Hence

\[
b_n
=
\frac{2(-1)^{n+1}}{n}.
\]

Therefore

\begin{answerbox}
\[
\boxed{
x
\sim
2\sum_{n=1}^{\infty}
\frac{(-1)^{n+1}}{n}\sin(nx),
\qquad
-\pi<x<\pi.
}
\]
\end{answerbox}
\end{workedexamplebox}

%%%%%%%%%%%%%%%%%%%%%%%%%%%%%%%%%%%%%%%%%%%%%%%%%%%%%%%%%%%%
\subsection{Half-Range Expansions}
\label{subsec:half-range-fourier}
%%%%%%%%%%%%%%%%%%%%%%%%%%%%%%%%%%%%%%%%%%%%%%%%%%%%%%%%%%%%

Suppose \(f\) is originally defined only on

\[
[0,L].
\]

One may extend \(f\) to \([-L,L]\) in different ways.

An even extension produces a cosine series:

\[
\boxed{
f(x)
\sim
\frac{a_0}{2}
+
\sum_{n=1}^{\infty}
a_n
\cos\left(\frac{n\pi x}{L}\right).
}
\]

An odd extension produces a sine series:

\[
\boxed{
f(x)
\sim
\sum_{n=1}^{\infty}
b_n
\sin\left(\frac{n\pi x}{L}\right).
}
\]

Half-range expansions are especially important in boundary-value
problems.

%%%%%%%%%%%%%%%%%%%%%%%%%%%%%%%%%%%%%%%%%%%%%%%%%%%%%%%%%%%%
\subsection{Complex Fourier Series}
\label{subsec:complex-fourier-series}
%%%%%%%%%%%%%%%%%%%%%%%%%%%%%%%%%%%%%%%%%%%%%%%%%%%%%%%%%%%%

Euler's formula gives

\[
e^{inx}
=
\cos(nx)+i\sin(nx).
\]

Therefore Fourier series can be written more compactly as

\[
\boxed{
f(x)
\sim
\sum_{n=-\infty}^{\infty}
c_ne^{inx}.
}
\]

The complex Fourier coefficients are

\[
\boxed{
c_n
=
\frac{1}{2\pi}
\int_{-\pi}^{\pi}
f(x)e^{-inx}\,dx.
}
\]

This formulation is usually more natural in advanced analysis.

%%%%%%%%%%%%%%%%%%%%%%%%%%%%%%%%%%%%%%%%%%%%%%%%%%%%%%%%%%%%
\subsection{Relationship Between Real and Complex Coefficients}
\label{subsec:real-complex-fourier-coefficients}
%%%%%%%%%%%%%%%%%%%%%%%%%%%%%%%%%%%%%%%%%%%%%%%%%%%%%%%%%%%%

For \(n>0\),

\[
\boxed{
c_n
=
\frac{a_n-ib_n}{2},
}
\]

and

\[
\boxed{
c_{-n}
=
\frac{a_n+ib_n}{2}.
}
\]

Also,

\[
\boxed{
c_0=\frac{a_0}{2}.
}
\]

If \(f\) is real-valued, then

\[
\boxed{
c_{-n}
=
\overline{c_n}.
}
\]

This conjugate symmetry ensures that the complex exponential expansion
represents a real-valued function.

%%%%%%%%%%%%%%%%%%%%%%%%%%%%%%%%%%%%%%%%%%%%%%%%%%%%%%%%%%%%
\subsection{Fourier Series as Hilbert-Space Expansions}
\label{subsec:fourier-hilbert-expansion}
%%%%%%%%%%%%%%%%%%%%%%%%%%%%%%%%%%%%%%%%%%%%%%%%%%%%%%%%%%%%

Consider

\[
L^2([-\pi,\pi])
\]

with inner product

\[
\boxed{
\langle f,g\rangle
=
\int_{-\pi}^{\pi}
f(x)\overline{g(x)}\,dx.
}
\]

Define

\[
\boxed{
e_n(x)
=
\frac{1}{\sqrt{2\pi}}e^{inx},
\qquad
n\in\Z.
}
\]

Then

\[
\boxed{
\langle e_n,e_m\rangle
=
\delta_{nm}.
}
\]

Thus

\[
\{e_n\}_{n\in\Z}
\]

is an orthonormal system.

In fact, it is an orthonormal basis of

\[
L^2([-\pi,\pi]).
\]

%%%%%%%%%%%%%%%%%%%%%%%%%%%%%%%%%%%%%%%%%%%%%%%%%%%%%%%%%%%%
\subsection{Fourier Coefficients as Coordinates}
\label{subsec:fourier-coordinates}
%%%%%%%%%%%%%%%%%%%%%%%%%%%%%%%%%%%%%%%%%%%%%%%%%%%%%%%%%%%%

With the normalized basis

\[
e_n(x)
=
\frac{1}{\sqrt{2\pi}}e^{inx},
\]

the Hilbert-space coordinate of \(f\) is

\[
\langle f,e_n\rangle
=
\frac{1}{\sqrt{2\pi}}
\int_{-\pi}^{\pi}
f(x)e^{-inx}\,dx.
\]

Since

\[
c_n
=
\frac{1}{2\pi}
\int_{-\pi}^{\pi}
f(x)e^{-inx}\,dx,
\]

we have

\[
\boxed{
\langle f,e_n\rangle
=
\sqrt{2\pi}\,c_n.
}
\]

Fourier coefficients are therefore infinite-dimensional coordinates.

%%%%%%%%%%%%%%%%%%%%%%%%%%%%%%%%%%%%%%%%%%%%%%%%%%%%%%%%%%%%
\subsection{Bessel's Inequality for Fourier Series}
\label{subsec:bessel-fourier}
%%%%%%%%%%%%%%%%%%%%%%%%%%%%%%%%%%%%%%%%%%%%%%%%%%%%%%%%%%%%

Bessel's inequality gives

\[
\sum_{n\in\Z}
|\langle f,e_n\rangle|^2
\leq
\|f\|_2^2.
\]

In terms of complex Fourier coefficients,

\[
\boxed{
2\pi
\sum_{n\in\Z}|c_n|^2
\leq
\int_{-\pi}^{\pi}|f(x)|^2\,dx.
}
\]

This gives quantitative control over Fourier coefficients.

%%%%%%%%%%%%%%%%%%%%%%%%%%%%%%%%%%%%%%%%%%%%%%%%%%%%%%%%%%%%
\subsection{Parseval's Identity for Fourier Series}
\label{subsec:parseval-fourier-series}
%%%%%%%%%%%%%%%%%%%%%%%%%%%%%%%%%%%%%%%%%%%%%%%%%%%%%%%%%%%%

Completeness of the trigonometric system strengthens Bessel's inequality
to equality.

\begin{theorem}[Parseval's Identity]
For \(f\in L^2([-\pi,\pi])\),

\[
\boxed{
\int_{-\pi}^{\pi}|f(x)|^2\,dx
=
2\pi
\sum_{n=-\infty}^{\infty}|c_n|^2.
}
\]
\end{theorem}

In real form,

\[
\boxed{
\frac1\pi
\int_{-\pi}^{\pi}|f(x)|^2\,dx
=
\frac{|a_0|^2}{2}
+
\sum_{n=1}^{\infty}
\left(
|a_n|^2+|b_n|^2
\right).
}
\]

Parseval's identity equates energy in the original function with energy
distributed among its frequencies.

%%%%%%%%%%%%%%%%%%%%%%%%%%%%%%%%%%%%%%%%%%%%%%%%%%%%%%%%%%%%
\subsection{Convergence of Fourier Series}
\label{subsec:convergence-fourier-series}
%%%%%%%%%%%%%%%%%%%%%%%%%%%%%%%%%%%%%%%%%%%%%%%%%%%%%%%%%%%%

The notation

\[
f(x)\sim\sum c_ne^{inx}
\]

does not automatically mean that the series converges pointwise to \(f\).

Several notions of convergence are relevant:

\[
\boxed{
\text{pointwise convergence},
}
\]

\[
\boxed{
\text{uniform convergence},
}
\]

and

\[
\boxed{
L^2\text{ convergence}.
}
\]

These are different concepts and require different hypotheses.

%%%%%%%%%%%%%%%%%%%%%%%%%%%%%%%%%%%%%%%%%%%%%%%%%%%%%%%%%%%%
\subsection{\(L^2\) Convergence}
\label{subsec:l2-fourier-convergence}
%%%%%%%%%%%%%%%%%%%%%%%%%%%%%%%%%%%%%%%%%%%%%%%%%%%%%%%%%%%%

For

\[
f\in L^2([-\pi,\pi]),
\]

let

\[
S_Nf
=
\sum_{n=-N}^{N}c_ne^{inx}.
\]

Completeness of the Fourier basis implies

\[
\boxed{
\|S_Nf-f\|_2\to0.
}
\]

Thus Fourier series converge to \(f\) in the \(L^2\) sense for every
\(L^2\) function.

This does not by itself imply pointwise convergence at every point.

%%%%%%%%%%%%%%%%%%%%%%%%%%%%%%%%%%%%%%%%%%%%%%%%%%%%%%%%%%%%
\subsection{Dirichlet Convergence Principle}
\label{subsec:dirichlet-convergence-fourier}
%%%%%%%%%%%%%%%%%%%%%%%%%%%%%%%%%%%%%%%%%%%%%%%%%%%%%%%%%%%%

For sufficiently regular periodic functions, Fourier series have a
classical pointwise convergence description.

At a point \(x\), under standard Dirichlet-type hypotheses,

\[
\boxed{
S_Nf(x)
\to
\frac{
f(x^-)+f(x^+)
}{2}.
}
\]

Therefore, at a point of continuity,

\[
\boxed{
S_Nf(x)\to f(x).
}
\]

At a jump discontinuity, the Fourier series converges to the midpoint of
the jump.

%%%%%%%%%%%%%%%%%%%%%%%%%%%%%%%%%%%%%%%%%%%%%%%%%%%%%%%%%%%%
\subsection{Dirichlet Kernel}
\label{subsec:dirichlet-kernel}
%%%%%%%%%%%%%%%%%%%%%%%%%%%%%%%%%%%%%%%%%%%%%%%%%%%%%%%%%%%%

The \(N\)th partial sum can be written as

\[
S_Nf(x)
=
\sum_{n=-N}^{N}c_ne^{inx}.
\]

Substituting the coefficient formula gives

\[
\boxed{
S_Nf(x)
=
\frac{1}{2\pi}
\int_{-\pi}^{\pi}
f(x-t)D_N(t)\,dt,
}
\]

where

\[
\boxed{
D_N(t)
=
\sum_{n=-N}^{N}e^{int}.
}
\]

The function \(D_N\) is the Dirichlet kernel.

For \(t\notin2\pi\Z\),

\[
\boxed{
D_N(t)
=
\frac{
\sin\left((N+\frac12)t\right)
}{
\sin(t/2)
}.
}
\]

%%%%%%%%%%%%%%%%%%%%%%%%%%%%%%%%%%%%%%%%%%%%%%%%%%%%%%%%%%%%
\subsection{Gibbs Phenomenon}
\label{subsec:gibbs-phenomenon}
%%%%%%%%%%%%%%%%%%%%%%%%%%%%%%%%%%%%%%%%%%%%%%%%%%%%%%%%%%%%

Near a jump discontinuity, Fourier partial sums develop oscillations and
overshoot the jump.

This is called the Gibbs phenomenon.

Increasing the number of Fourier terms makes the oscillatory region
narrower, but the relative overshoot does not vanish.

Thus

\[
\boxed{
\text{more terms}
\not\Rightarrow
\text{uniform convergence near a jump}.
}
\]

The phenomenon is important in signal reconstruction and numerical
approximation.

%%%%%%%%%%%%%%%%%%%%%%%%%%%%%%%%%%%%%%%%%%%%%%%%%%%%%%%%%%%%
\subsection{Fej\'er Means}
\label{subsec:fejer-means}
%%%%%%%%%%%%%%%%%%%%%%%%%%%%%%%%%%%%%%%%%%%%%%%%%%%%%%%%%%%%

Instead of using a single partial sum, one may average the first several
partial sums.

Define

\[
\boxed{
\sigma_Nf
=
\frac{1}{N+1}
\sum_{k=0}^{N}S_kf.
}
\]

The lowercase Greek letter \(\sigma\) is \emph{sigma}.

These are called Fej\'er means.

%%%%%%%%%%%%%%%%%%%%%%%%%%%%%%%%%%%%%%%%%%%%%%%%%%%%%%%%%%%%
\subsection{Fej\'er's Theorem}
\label{subsec:fejer-theorem}
%%%%%%%%%%%%%%%%%%%%%%%%%%%%%%%%%%%%%%%%%%%%%%%%%%%%%%%%%%%%

\begin{theorem}[Fej\'er's Theorem]
If \(f\) is continuous and \(2\pi\)-periodic, then its Fej\'er means
converge uniformly to \(f\):

\[
\boxed{
\|\sigma_Nf-f\|_\infty\to0.
}
\]
\end{theorem}

Fej\'er summation improves convergence by replacing the oscillatory
Dirichlet kernel with a positive averaging kernel.

%%%%%%%%%%%%%%%%%%%%%%%%%%%%%%%%%%%%%%%%%%%%%%%%%%%%%%%%%%%%
\subsection{Convolution on the Circle}
\label{subsec:periodic-convolution}
%%%%%%%%%%%%%%%%%%%%%%%%%%%%%%%%%%%%%%%%%%%%%%%%%%%%%%%%%%%%

For \(2\pi\)-periodic functions, define convolution by

\[
\boxed{
(f*g)(x)
=
\frac{1}{2\pi}
\int_{-\pi}^{\pi}
f(x-t)g(t)\,dt.
}
\]

Then the Fourier coefficients satisfy

\[
\boxed{
\widehat{f*g}(n)
=
\widehat{f}(n)\widehat{g}(n),
}
\]

under this normalization.

Thus convolution in the original variable becomes multiplication in
frequency.

%%%%%%%%%%%%%%%%%%%%%%%%%%%%%%%%%%%%%%%%%%%%%%%%%%%%%%%%%%%%
\subsection{Approximate Identities}
\label{subsec:approximate-identities}
%%%%%%%%%%%%%%%%%%%%%%%%%%%%%%%%%%%%%%%%%%%%%%%%%%%%%%%%%%%%

An approximate identity is a family of kernels that increasingly
concentrates near the origin while preserving total mass.

Schematically,

\[
\boxed{
K_\varepsilon*f
\longrightarrow
f
}
\]

as

\[
\varepsilon\to0.
\]

Approximate identities provide a general framework for smoothing and
reconstruction.

Fej\'er kernels are a periodic example.

%%%%%%%%%%%%%%%%%%%%%%%%%%%%%%%%%%%%%%%%%%%%%%%%%%%%%%%%%%%%
\subsection{From Fourier Series to the Fourier Transform}
\label{subsec:series-to-transform}
%%%%%%%%%%%%%%%%%%%%%%%%%%%%%%%%%%%%%%%%%%%%%%%%%%%%%%%%%%%%

Fourier series describe periodic functions using discrete frequencies

\[
n\in\Z.
\]

For nonperiodic functions on the entire real line, the frequency variable
becomes continuous.

This leads to the Fourier transform.

Schematically,

\[
\boxed{
\text{Fourier series}
\quad
\sum_{n\in\Z}
}
\]

becomes

\[
\boxed{
\text{Fourier transform}
\quad
\int_{\R}.
}
\]

%%%%%%%%%%%%%%%%%%%%%%%%%%%%%%%%%%%%%%%%%%%%%%%%%%%%%%%%%%%%
\subsection{Fourier Transform}
\label{subsec:fourier-transform}
%%%%%%%%%%%%%%%%%%%%%%%%%%%%%%%%%%%%%%%%%%%%%%%%%%%%%%%%%%%%

We adopt the convention

\[
\boxed{
\widehat{f}(\xi)
=
\int_{-\infty}^{\infty}
f(x)e^{-2\pi i x\xi}\,dx.
}
\]

The lowercase Greek letter \(\xi\) is \emph{xi} and represents frequency.

The function

\[
\widehat{f}
\]

is called the Fourier transform of \(f\).

The hat notation is read ``\(f\)-hat.''

%%%%%%%%%%%%%%%%%%%%%%%%%%%%%%%%%%%%%%%%%%%%%%%%%%%%%%%%%%%%
\subsection{Inverse Fourier Transform}
\label{subsec:inverse-fourier-transform}
%%%%%%%%%%%%%%%%%%%%%%%%%%%%%%%%%%%%%%%%%%%%%%%%%%%%%%%%%%%%

Under suitable hypotheses,

\[
\boxed{
f(x)
=
\int_{-\infty}^{\infty}
\widehat{f}(\xi)
e^{2\pi i x\xi}
\,d\xi.
}
\]

Thus the Fourier transform and inverse transform form a dual pair:

\[
\boxed{
f(x)
\longleftrightarrow
\widehat{f}(\xi).
}
\]

The exact constants depend on the Fourier-transform convention being used.

%%%%%%%%%%%%%%%%%%%%%%%%%%%%%%%%%%%%%%%%%%%%%%%%%%%%%%%%%%%%
\subsection{Basic Transform Pair}
\label{subsec:basic-transform-pair}
%%%%%%%%%%%%%%%%%%%%%%%%%%%%%%%%%%%%%%%%%%%%%%%%%%%%%%%%%%%%

The Gaussian plays a distinguished role in Fourier analysis.

Under the convention used here,

\[
\boxed{
f(x)=e^{-\pi x^2}
}
\]

satisfies

\[
\boxed{
\widehat{f}(\xi)
=
e^{-\pi\xi^2}.
}
\]

Thus this Gaussian is its own Fourier transform.

%%%%%%%%%%%%%%%%%%%%%%%%%%%%%%%%%%%%%%%%%%%%%%%%%%%%%%%%%%%%
\subsection{Linearity}
\label{subsec:fourier-linearity}
%%%%%%%%%%%%%%%%%%%%%%%%%%%%%%%%%%%%%%%%%%%%%%%%%%%%%%%%%%%%

The Fourier transform is linear:

\[
\boxed{
\widehat{af+bg}
=
a\widehat{f}
+
b\widehat{g}.
}
\]

This follows immediately from linearity of integration.

%%%%%%%%%%%%%%%%%%%%%%%%%%%%%%%%%%%%%%%%%%%%%%%%%%%%%%%%%%%%
\subsection{Translation}
\label{subsec:fourier-translation}
%%%%%%%%%%%%%%%%%%%%%%%%%%%%%%%%%%%%%%%%%%%%%%%%%%%%%%%%%%%%

If

\[
g(x)=f(x-a),
\]

then

\[
\boxed{
\widehat{g}(\xi)
=
e^{-2\pi i a\xi}
\widehat{f}(\xi).
}
\]

Thus translation in the original variable corresponds to multiplication
by a phase factor in frequency.

%%%%%%%%%%%%%%%%%%%%%%%%%%%%%%%%%%%%%%%%%%%%%%%%%%%%%%%%%%%%
\subsection{Modulation}
\label{subsec:fourier-modulation}
%%%%%%%%%%%%%%%%%%%%%%%%%%%%%%%%%%%%%%%%%%%%%%%%%%%%%%%%%%%%

If

\[
g(x)
=
e^{2\pi i ax}f(x),
\]

then

\[
\boxed{
\widehat{g}(\xi)
=
\widehat{f}(\xi-a).
}
\]

Thus multiplication by a complex oscillation shifts the frequency
spectrum.

This operation is called modulation.

%%%%%%%%%%%%%%%%%%%%%%%%%%%%%%%%%%%%%%%%%%%%%%%%%%%%%%%%%%%%
\subsection{Scaling}
\label{subsec:fourier-scaling}
%%%%%%%%%%%%%%%%%%%%%%%%%%%%%%%%%%%%%%%%%%%%%%%%%%%%%%%%%%%%

If

\[
g(x)=f(ax),
\qquad
a\neq0,
\]

then

\[
\boxed{
\widehat{g}(\xi)
=
\frac{1}{|a|}
\widehat{f}\left(\frac{\xi}{a}\right).
}
\]

Therefore compression in the original variable produces expansion in
frequency, and vice versa.

This is one manifestation of the uncertainty principle.

%%%%%%%%%%%%%%%%%%%%%%%%%%%%%%%%%%%%%%%%%%%%%%%%%%%%%%%%%%%%
\subsection{Differentiation}
\label{subsec:fourier-differentiation}
%%%%%%%%%%%%%%%%%%%%%%%%%%%%%%%%%%%%%%%%%%%%%%%%%%%%%%%%%%%%

Suppose \(f\) and its derivative have sufficient decay and regularity.

Integration by parts gives

\[
\boxed{
\widehat{f'}(\xi)
=
2\pi i\xi\,\widehat{f}(\xi).
}
\]

More generally,

\[
\boxed{
\widehat{f^{(n)}}(\xi)
=
(2\pi i\xi)^n
\widehat{f}(\xi).
}
\]

Thus differentiation becomes multiplication by a polynomial in frequency.

This is one of the most important reasons Fourier methods are effective
for differential equations.

%%%%%%%%%%%%%%%%%%%%%%%%%%%%%%%%%%%%%%%%%%%%%%%%%%%%%%%%%%%%
\subsection{Multiplication by the Independent Variable}
\label{subsec:fourier-multiplication-x}
%%%%%%%%%%%%%%%%%%%%%%%%%%%%%%%%%%%%%%%%%%%%%%%%%%%%%%%%%%%%

Under suitable hypotheses,

\[
\boxed{
\widehat{xf(x)}(\xi)
=
-\frac{1}{2\pi i}
\frac{d}{d\xi}
\widehat{f}(\xi).
}
\]

Thus multiplication by \(x\) in physical space corresponds to
differentiation in frequency space.

Together with the derivative rule, this reveals a fundamental duality:

\[
\boxed{
\text{differentiation}
\longleftrightarrow
\text{multiplication}.
}
\]

%%%%%%%%%%%%%%%%%%%%%%%%%%%%%%%%%%%%%%%%%%%%%%%%%%%%%%%%%%%%
\subsection{Convolution on the Real Line}
\label{subsec:convolution-real-line}
%%%%%%%%%%%%%%%%%%%%%%%%%%%%%%%%%%%%%%%%%%%%%%%%%%%%%%%%%%%%

For suitable functions \(f\) and \(g\), define

\[
\boxed{
(f*g)(x)
=
\int_{-\infty}^{\infty}
f(x-y)g(y)\,dy.
}
\]

Convolution combines the values of two functions through translation and
averaging.

It appears throughout analysis, probability, PDEs, and signal processing.

%%%%%%%%%%%%%%%%%%%%%%%%%%%%%%%%%%%%%%%%%%%%%%%%%%%%%%%%%%%%
\subsection{Convolution Theorem}
\label{subsec:convolution-theorem-fourier}
%%%%%%%%%%%%%%%%%%%%%%%%%%%%%%%%%%%%%%%%%%%%%%%%%%%%%%%%%%%%

\begin{theorem}[Convolution Theorem]
Under suitable integrability assumptions,

\[
\boxed{
\widehat{f*g}
=
\widehat{f}\,\widehat{g}.
}
\]
\end{theorem}

Also,

\[
\boxed{
\widehat{fg}
=
\widehat{f}*\widehat{g}
}
\]

when the expressions are defined appropriately.

Thus Fourier transformation exchanges multiplication and convolution.

%%%%%%%%%%%%%%%%%%%%%%%%%%%%%%%%%%%%%%%%%%%%%%%%%%%%%%%%%%%%
\subsection{Young's Convolution Inequality}
\label{subsec:young-convolution}
%%%%%%%%%%%%%%%%%%%%%%%%%%%%%%%%%%%%%%%%%%%%%%%%%%%%%%%%%%%%

One useful estimate is

\[
\boxed{
\|f*g\|_r
\leq
\|f\|_p\|g\|_q,
}
\]

under the exponent relation

\[
\boxed{
1+\frac1r
=
\frac1p+\frac1q,
}
\]

for appropriate \(p,q,r\).

A particularly important case is

\[
\boxed{
\|f*g\|_p
\leq
\|f\|_p\|g\|_1.
}
\]

This estimate is central to smoothing and approximation arguments.

%%%%%%%%%%%%%%%%%%%%%%%%%%%%%%%%%%%%%%%%%%%%%%%%%%%%%%%%%%%%
\subsection{Riemann--Lebesgue Lemma}
\label{subsec:riemann-lebesgue}
%%%%%%%%%%%%%%%%%%%%%%%%%%%%%%%%%%%%%%%%%%%%%%%%%%%%%%%%%%%%

\begin{theorem}[Riemann--Lebesgue Lemma]
If

\[
f\in L^1(\R),
\]

then

\[
\boxed{
\widehat{f}(\xi)\to0
\qquad
\text{as }
|\xi|\to\infty.
}
\]
\end{theorem}

Thus integrable functions have Fourier transforms that vanish at high
frequency.

The theorem does not state that every continuous function vanishing at
infinity is the Fourier transform of an \(L^1\) function.

%%%%%%%%%%%%%%%%%%%%%%%%%%%%%%%%%%%%%%%%%%%%%%%%%%%%%%%%%%%%
\subsection{Plancherel's Theorem}
\label{subsec:plancherel-theorem}
%%%%%%%%%%%%%%%%%%%%%%%%%%%%%%%%%%%%%%%%%%%%%%%%%%%%%%%%%%%%

The Fourier transform extends from suitable integrable functions to a
unitary operator on \(L^2(\R)\).

\begin{theorem}[Plancherel's Theorem]
For \(f\in L^2(\R)\),

\[
\boxed{
\|\widehat{f}\|_2
=
\|f\|_2.
}
\]
\end{theorem}

Equivalently,

\[
\boxed{
\int_{-\infty}^{\infty}
|\widehat{f}(\xi)|^2\,d\xi
=
\int_{-\infty}^{\infty}
|f(x)|^2\,dx.
}
\]

This is the continuous-frequency analogue of Parseval's identity.

%%%%%%%%%%%%%%%%%%%%%%%%%%%%%%%%%%%%%%%%%%%%%%%%%%%%%%%%%%%%
\subsection{Fourier Transform as a Unitary Operator}
\label{subsec:fourier-unitary}
%%%%%%%%%%%%%%%%%%%%%%%%%%%%%%%%%%%%%%%%%%%%%%%%%%%%%%%%%%%%

Plancherel's Theorem allows the Fourier transform to be interpreted as

\[
\boxed{
\mathcal{F}:L^2(\R)\to L^2(\R),
}
\]

where

\[
\mathcal{F}f=\widehat{f}.
\]

Under the chosen normalization,

\[
\boxed{
\|\mathcal{F}f\|_2=\|f\|_2.
}
\]

Thus the Fourier transform preserves Hilbert-space energy.

This is a central connection between Fourier analysis and functional
analysis.

%%%%%%%%%%%%%%%%%%%%%%%%%%%%%%%%%%%%%%%%%%%%%%%%%%%%%%%%%%%%
\subsection{Fourier Inversion}
\label{subsec:fourier-inversion}
%%%%%%%%%%%%%%%%%%%%%%%%%%%%%%%%%%%%%%%%%%%%%%%%%%%%%%%%%%%%

Under appropriate hypotheses, Fourier transformation can be reversed.

A standard inversion statement is

\[
\boxed{
f(x)
=
\int_{\R}
\widehat{f}(\xi)
e^{2\pi i x\xi}
\,d\xi.
}
\]

For example, if

\[
f,\widehat{f}\in L^1(\R),
\]

then inversion holds at suitable points, and in particular at continuity
points under standard assumptions.

More general versions hold in \(L^2\) and distributional settings.

%%%%%%%%%%%%%%%%%%%%%%%%%%%%%%%%%%%%%%%%%%%%%%%%%%%%%%%%%%%%
\subsection{Schwartz Functions}
\label{subsec:schwartz-functions}
%%%%%%%%%%%%%%%%%%%%%%%%%%%%%%%%%%%%%%%%%%%%%%%%%%%%%%%%%%%%

Fourier transformation behaves especially well on rapidly decreasing
smooth functions.

\begin{definitionbox}{Schwartz Space}
The Schwartz space

\[
\boxed{
\mathcal{S}(\R)
}
\]

consists of smooth functions \(f\) such that

\[
\boxed{
\sup_{x\in\R}
|x^m f^{(n)}(x)|
<\infty
}
\]

for every pair of nonnegative integers \(m,n\).
\end{definitionbox}

Schwartz functions and all their derivatives decay faster than every
inverse polynomial.

The Fourier transform maps

\[
\boxed{
\mathcal{S}(\R)
\to
\mathcal{S}(\R).
}
\]

%%%%%%%%%%%%%%%%%%%%%%%%%%%%%%%%%%%%%%%%%%%%%%%%%%%%%%%%%%%%
\subsection{Fourier Transform of the Gaussian}
\label{subsec:fourier-gaussian}
%%%%%%%%%%%%%%%%%%%%%%%%%%%%%%%%%%%%%%%%%%%%%%%%%%%%%%%%%%%%

Let

\[
f(x)=e^{-\pi x^2}.
\]

Then

\[
\boxed{
\widehat{f}(\xi)=e^{-\pi\xi^2}.
}
\]

More generally, for \(a>0\),

\[
\boxed{
\mathcal{F}
\left(
e^{-\pi a x^2}
\right)(\xi)
=
a^{-1/2}
e^{-\pi\xi^2/a}.
}
\]

The Gaussian is therefore naturally compatible with Fourier
transformation.

This fact is fundamental in probability, PDEs, quantum mechanics, and
harmonic analysis.

%%%%%%%%%%%%%%%%%%%%%%%%%%%%%%%%%%%%%%%%%%%%%%%%%%%%%%%%%%%%
\subsection{Uncertainty Principle}
\label{subsec:uncertainty-principle-fourier}
%%%%%%%%%%%%%%%%%%%%%%%%%%%%%%%%%%%%%%%%%%%%%%%%%%%%%%%%%%%%

A function and its Fourier transform cannot both be arbitrarily
concentrated.

One quantitative form states that, under suitable normalization and
regularity assumptions,

\[
\boxed{
\left(
\int_{\R}
(x-a)^2|f(x)|^2\,dx
\right)
\left(
\int_{\R}
(\xi-b)^2|\widehat{f}(\xi)|^2\,d\xi
\right)
\geq
\frac{\|f\|_2^4}{16\pi^2}.
}
\]

Here \(a\) and \(b\) may be chosen as appropriate centers.

The Gaussian is the extremizing shape for the classical uncertainty
inequality.

Conceptually,

\[
\boxed{
\text{localization in one domain}
\Longrightarrow
\text{dispersion in the other}.
}
\]

%%%%%%%%%%%%%%%%%%%%%%%%%%%%%%%%%%%%%%%%%%%%%%%%%%%%%%%%%%%%
\subsection{Poisson Summation Formula}
\label{subsec:poisson-summation}
%%%%%%%%%%%%%%%%%%%%%%%%%%%%%%%%%%%%%%%%%%%%%%%%%%%%%%%%%%%%

Under suitable decay and regularity assumptions,

\[
\boxed{
\sum_{n\in\Z}f(n)
=
\sum_{k\in\Z}\widehat{f}(k).
}
\]

This is the Poisson summation formula.

It connects

\[
\boxed{
\text{sampling in physical space}
}
\]

with

\[
\boxed{
\text{periodization in frequency space}.
}
\]

It has important applications in number theory, signal processing,
theta functions, and lattice sums.

%%%%%%%%%%%%%%%%%%%%%%%%%%%%%%%%%%%%%%%%%%%%%%%%%%%%%%%%%%%%
\subsection{Fourier Methods for Differential Equations}
\label{subsec:fourier-differential-equations}
%%%%%%%%%%%%%%%%%%%%%%%%%%%%%%%%%%%%%%%%%%%%%%%%%%%%%%%%%%%%

Suppose a differential equation has the schematic form

\[
P(D)u=f,
\]

where \(P(D)\) is a constant-coefficient differential operator.

Fourier transformation converts derivatives into multiplication:

\[
\widehat{u^{(n)}}(\xi)
=
(2\pi i\xi)^n\widehat{u}(\xi).
\]

Therefore the differential equation becomes an algebraic equation

\[
\boxed{
P(2\pi i\xi)\widehat{u}(\xi)
=
\widehat{f}(\xi).
}
\]

Formally,

\[
\boxed{
\widehat{u}(\xi)
=
\frac{\widehat{f}(\xi)}
{P(2\pi i\xi)}.
}
\]

The solution is then recovered using the inverse Fourier transform.

%%%%%%%%%%%%%%%%%%%%%%%%%%%%%%%%%%%%%%%%%%%%%%%%%%%%%%%%%%%%
\subsection{The Heat Equation}
\label{subsec:fourier-heat-equation}
%%%%%%%%%%%%%%%%%%%%%%%%%%%%%%%%%%%%%%%%%%%%%%%%%%%%%%%%%%%%

Consider the heat equation on the real line,

\[
\boxed{
u_t
=
\kappa u_{xx},
}
\]

where \(\kappa>0\) is the diffusion coefficient.

Taking the Fourier transform in \(x\) gives

\[
\frac{\partial}{\partial t}
\widehat{u}(\xi,t)
=
-4\pi^2\kappa\xi^2
\widehat{u}(\xi,t).
\]

For each fixed \(\xi\), this is an ordinary differential equation.

Hence

\[
\boxed{
\widehat{u}(\xi,t)
=
e^{-4\pi^2\kappa\xi^2t}
\widehat{u}_0(\xi).
}
\]

Thus high frequencies decay faster than low frequencies.

This explains mathematically why diffusion smooths functions.

%%%%%%%%%%%%%%%%%%%%%%%%%%%%%%%%%%%%%%%%%%%%%%%%%%%%%%%%%%%%
\subsection{The Wave Equation}
\label{subsec:fourier-wave-equation}
%%%%%%%%%%%%%%%%%%%%%%%%%%%%%%%%%%%%%%%%%%%%%%%%%%%%%%%%%%%%

Consider

\[
\boxed{
u_{tt}
=
c^2u_{xx}.
}
\]

Taking the Fourier transform in the spatial variable gives

\[
\widehat{u}_{tt}
=
-4\pi^2c^2\xi^2\widehat{u}.
\]

Thus each frequency satisfies a harmonic oscillator equation.

Fourier analysis therefore decomposes the wave equation into independent
oscillatory modes.

%%%%%%%%%%%%%%%%%%%%%%%%%%%%%%%%%%%%%%%%%%%%%%%%%%%%%%%%%%%%
\subsection{Frequency Filtering}
\label{subsec:frequency-filtering}
%%%%%%%%%%%%%%%%%%%%%%%%%%%%%%%%%%%%%%%%%%%%%%%%%%%%%%%%%%%%

Suppose a system acts on a signal by convolution:

\[
y=h*x.
\]

Taking Fourier transforms gives

\[
\boxed{
\widehat{y}
=
\widehat{h}\,\widehat{x}.
}
\]

Thus a convolution operator becomes multiplication by the frequency
response

\[
\widehat{h}.
\]

Filters can therefore amplify, suppress, or remove selected frequency
ranges.

This principle underlies much of signal processing.

%%%%%%%%%%%%%%%%%%%%%%%%%%%%%%%%%%%%%%%%%%%%%%%%%%%%%%%%%%%%
\subsection{Discrete Fourier Transform Preview}
\label{subsec:dft-preview}
%%%%%%%%%%%%%%%%%%%%%%%%%%%%%%%%%%%%%%%%%%%%%%%%%%%%%%%%%%%%

For sampled data

\[
x_0,x_1,\ldots,x_{N-1},
\]

the discrete Fourier transform is

\[
\boxed{
X_k
=
\sum_{n=0}^{N-1}
x_n
e^{-2\pi i kn/N},
\qquad
k=0,\ldots,N-1.
}
\]

The inverse discrete Fourier transform is

\[
\boxed{
x_n
=
\frac1N
\sum_{k=0}^{N-1}
X_k
e^{2\pi i kn/N}.
}
\]

The DFT converts a finite data vector from a sample representation to a
frequency representation.

Its computational implementation through the Fast Fourier Transform is
developed more naturally in Numerical and Computational Mathematics.

%%%%%%%%%%%%%%%%%%%%%%%%%%%%%%%%%%%%%%%%%%%%%%%%%%%%%%%%%%%%
\subsection{Fourier Analysis and Probability}
\label{subsec:fourier-probability}
%%%%%%%%%%%%%%%%%%%%%%%%%%%%%%%%%%%%%%%%%%%%%%%%%%%%%%%%%%%%

For a probability density \(f\), the Fourier transform is closely related
to the characteristic function of a random variable.

Under a common probability convention,

\[
\boxed{
\varphi_X(t)
=
\mathbb{E}[e^{itX}].
}
\]

This is essentially a Fourier transform of the probability distribution.

Independence and addition interact naturally with multiplication:

\[
\boxed{
\varphi_{X+Y}(t)
=
\varphi_X(t)\varphi_Y(t)
}
\]

when \(X\) and \(Y\) are independent.

Fourier methods therefore play a central role in limit theorems and
distribution theory.

%%%%%%%%%%%%%%%%%%%%%%%%%%%%%%%%%%%%%%%%%%%%%%%%%%%%%%%%%%%%
\subsection{Fourier Analysis and Complex Analysis}
\label{subsec:fourier-complex-analysis}
%%%%%%%%%%%%%%%%%%%%%%%%%%%%%%%%%%%%%%%%%%%%%%%%%%%%%%%%%%%%

Complex exponentials

\[
e^{iz}
\]

provide the natural algebraic language of Fourier analysis.

Contour integration can also be used to evaluate Fourier-type integrals.

Analytic continuation and growth estimates connect Fourier transforms to
complex-variable theory.

These ideas become especially important in advanced harmonic analysis and
the theory of distributions.

%%%%%%%%%%%%%%%%%%%%%%%%%%%%%%%%%%%%%%%%%%%%%%%%%%%%%%%%%%%%
\subsection{Fourier Analysis and Functional Analysis}
\label{subsec:fourier-functional-analysis}
%%%%%%%%%%%%%%%%%%%%%%%%%%%%%%%%%%%%%%%%%%%%%%%%%%%%%%%%%%%%

Fourier analysis can be interpreted as a change of coordinates in
function space.

For Fourier series,

\[
\boxed{
L^2([-\pi,\pi])
\longleftrightarrow
\ell^2(\Z).
}
\]

For the Fourier transform,

\[
\boxed{
L^2(\R)
\longrightarrow
L^2(\R)
}
\]

is unitary under the chosen normalization.

This makes Fourier analysis one of the most important concrete
realizations of Hilbert-space theory.

%%%%%%%%%%%%%%%%%%%%%%%%%%%%%%%%%%%%%%%%%%%%%%%%%%%%%%%%%%%%
\subsection{Common Errors}
\label{subsec:fourier-analysis-common-errors}
%%%%%%%%%%%%%%%%%%%%%%%%%%%%%%%%%%%%%%%%%%%%%%%%%%%%%%%%%%%%

\subsubsection{Assuming a Fourier Series Always Converges Pointwise}

An \(L^2\) Fourier series converges in \(L^2\), but this does not
automatically imply pointwise convergence everywhere.

%%%%%%%%%%%%%%%%%%%%%%%%%%%%%%%%%%%%%%%%%%%%%%%%%%%%%%%%%%%%
\subsubsection{Ignoring the Period}

The coefficient formulas depend on the interval and period.

For period \(2L\), the frequencies are

\[
\frac{n\pi}{L},
\]

not simply \(n\).

%%%%%%%%%%%%%%%%%%%%%%%%%%%%%%%%%%%%%%%%%%%%%%%%%%%%%%%%%%%%
\subsubsection{Missing Symmetry}

If \(f\) is even, all sine coefficients vanish.

If \(f\) is odd, all cosine coefficients and \(a_0\) vanish.

Checking symmetry before integrating can greatly simplify calculations.

%%%%%%%%%%%%%%%%%%%%%%%%%%%%%%%%%%%%%%%%%%%%%%%%%%%%%%%%%%%%
\subsubsection{Forgetting the Constant Term}

The real Fourier series contains

\[
\boxed{
\frac{a_0}{2},
}
\]

not \(a_0\).

%%%%%%%%%%%%%%%%%%%%%%%%%%%%%%%%%%%%%%%%%%%%%%%%%%%%%%%%%%%%
\subsubsection{Confusing Fourier Series with Fourier Transforms}

Fourier series use discrete frequencies for periodic functions.

Fourier transforms use a continuous frequency variable for functions on
the real line.

%%%%%%%%%%%%%%%%%%%%%%%%%%%%%%%%%%%%%%%%%%%%%%%%%%%%%%%%%%%%
\subsubsection{Mixing Fourier Transform Conventions}

Different texts distribute factors of

\[
2\pi
\]

differently between the transform and inverse transform.

All formulas must be consistent with the chosen convention.

%%%%%%%%%%%%%%%%%%%%%%%%%%%%%%%%%%%%%%%%%%%%%%%%%%%%%%%%%%%%
\subsubsection{Forgetting the Scaling Factor}

If

\[
g(x)=f(ax),
\]

then

\[
\widehat{g}(\xi)
=
\frac1{|a|}
\widehat{f}\left(\frac{\xi}{a}\right).
\]

The factor \(1/|a|\) is essential.

%%%%%%%%%%%%%%%%%%%%%%%%%%%%%%%%%%%%%%%%%%%%%%%%%%%%%%%%%%%%
\subsubsection{Confusing Multiplication and Convolution}

Under Fourier transformation,

\[
\boxed{
\text{convolution}
\longleftrightarrow
\text{multiplication}.
}
\]

These operations should not be interchanged without transforming.

%%%%%%%%%%%%%%%%%%%%%%%%%%%%%%%%%%%%%%%%%%%%%%%%%%%%%%%%%%%%
\subsubsection{Ignoring Regularity and Integrability Hypotheses}

Differentiation rules, inversion formulas, and interchange of integrals
require appropriate assumptions.

%%%%%%%%%%%%%%%%%%%%%%%%%%%%%%%%%%%%%%%%%%%%%%%%%%%%%%%%%%%%
\subsubsection{Misinterpreting Gibbs Oscillations}

Adding more terms narrows the region of oscillation near a jump but does
not eliminate the characteristic relative overshoot.

%%%%%%%%%%%%%%%%%%%%%%%%%%%%%%%%%%%%%%%%%%%%%%%%%%%%%%%%%%%%
\subsubsection{Assuming Pointwise Equality from \(L^2\) Equality}

Elements of \(L^2\) are equivalence classes modulo equality almost
everywhere.

An \(L^2\) identity need not be a pointwise identity at every location.

%%%%%%%%%%%%%%%%%%%%%%%%%%%%%%%%%%%%%%%%%%%%%%%%%%%%%%%%%%%%
\subsection{Applications}
\label{subsec:fourier-analysis-applications}
%%%%%%%%%%%%%%%%%%%%%%%%%%%%%%%%%%%%%%%%%%%%%%%%%%%%%%%%%%%%

\begin{applicationbox}
\textbf{Differential Equations.}
Fourier methods convert constant-coefficient differential operators into
multiplication operators in frequency space.

\medskip

\textbf{Heat and Diffusion.}
The heat equation suppresses high-frequency components rapidly, explaining
the smoothing effect of diffusion.

\medskip

\textbf{Wave Propagation.}
Fourier decomposition separates waves into independent frequency modes.

\medskip

\textbf{Signal Processing.}
Signals are decomposed into frequency components for filtering,
compression, reconstruction, and noise reduction.

\medskip

\textbf{Image Processing.}
Two-dimensional Fourier transforms are used for filtering, restoration,
frequency analysis, and compression.

\medskip

\textbf{Communications.}
Modulation, bandwidth, spectral density, and frequency filtering are
naturally expressed using Fourier methods.

\medskip

\textbf{Probability.}
Characteristic functions transform convolution of independent
distributions into multiplication.

\medskip

\textbf{Quantum Mechanics.}
Position and momentum representations are related by Fourier
transformation.

\medskip

\textbf{Numerical Computation.}
The discrete Fourier transform and FFT provide efficient algorithms for
large-scale frequency analysis.

\medskip

\textbf{Number Theory.}
Poisson summation and exponential sums connect Fourier analysis with
lattices, theta functions, and analytic number theory.
\end{applicationbox}

%%%%%%%%%%%%%%%%%%%%%%%%%%%%%%%%%%%%%%%%%%%%%%%%%%%%%%%%%%%%
\subsection{Exercises}
\label{subsec:fourier-analysis-exercises}
%%%%%%%%%%%%%%%%%%%%%%%%%%%%%%%%%%%%%%%%%%%%%%%%%%%%%%%%%%%%

\begin{exercise}[1]
Determine whether

\[
f(x)=\cos(3x)+2\sin(5x)
\]

is periodic and find a period.
\end{exercise}

\begin{exercise}[1]
Verify that

\[
\int_{-\pi}^{\pi}
\sin(mx)\cos(nx)\,dx=0
\]

for positive integers \(m,n\).
\end{exercise}

\begin{exercise}[1]
Determine whether

\[
f(x)=x^2
\]

is even, odd, or neither on \([-\pi,\pi]\).
\end{exercise}

\begin{exercise}[1]
Determine whether

\[
f(x)=x^3
\]

is even, odd, or neither.
\end{exercise}

\begin{exercise}[2]
Compute the Fourier series of

\[
f(x)=x
\]

on

\[
(-\pi,\pi).
\]
\end{exercise}

\begin{exercise}[2]
Compute the Fourier cosine series of

\[
f(x)=x
\]

on

\[
(0,\pi).
\]
\end{exercise}

\begin{exercise}[2]
Compute the Fourier sine series of

\[
f(x)=x
\]

on

\[
(0,\pi).
\]
\end{exercise}

\begin{exercise}[2]
Find the Fourier series of

\[
f(x)=x^2
\]

on

\[
(-\pi,\pi).
\]
\end{exercise}

\begin{exercise}[2]
Show that if \(f\) is real-valued, then its complex Fourier coefficients
satisfy

\[
c_{-n}=\overline{c_n}.
\]
\end{exercise}

\begin{exercise}[2]
Derive the relationship

\[
c_n=\frac{a_n-ib_n}{2}
\]

for \(n>0\).
\end{exercise}

\begin{exercise}[2]
Use the geometric-series identity to derive

\[
D_N(t)
=
\frac{
\sin\left((N+\frac12)t\right)
}{
\sin(t/2)
}.
\]
\end{exercise}

\begin{exercise}[2]
Explain what the Fourier series converges to at a jump discontinuity under
the standard Dirichlet hypotheses.
\end{exercise}

\begin{exercise}[2]
Use Parseval's identity and the Fourier series of \(x\) to derive a
classical infinite-series identity.
\end{exercise}

\begin{exercise}[3]
Prove that convolution on the circle is commutative.
\end{exercise}

\begin{exercise}[3]
Prove the convolution theorem for periodic Fourier coefficients.
\end{exercise}

\begin{exercise}[3]
Show that

\[
\widehat{f(\cdot-a)}(\xi)
=
e^{-2\pi ia\xi}\widehat{f}(\xi).
\]
\end{exercise}

\begin{exercise}[3]
Prove the modulation identity

\[
\mathcal{F}
\left(
e^{2\pi iax}f(x)
\right)(\xi)
=
\widehat{f}(\xi-a).
\]
\end{exercise}

\begin{exercise}[3]
Derive the Fourier scaling formula.
\end{exercise}

\begin{exercise}[3]
Assuming sufficient decay, prove

\[
\widehat{f'}(\xi)
=
2\pi i\xi\widehat{f}(\xi).
\]
\end{exercise}

\begin{exercise}[3]
Derive

\[
\widehat{xf(x)}(\xi)
=
-\frac{1}{2\pi i}
\widehat{f}'(\xi).
\]
\end{exercise}

\begin{exercise}[3]
Prove that convolution on \(\R\) is commutative whenever the integrals
exist.
\end{exercise}

\begin{exercise}[3]
Prove the Fourier convolution theorem under suitable \(L^1\) hypotheses.
\end{exercise}

\begin{exercise}[3]
Verify directly that

\[
e^{-\pi x^2}
\]

is its own Fourier transform.
\end{exercise}

\begin{exercise}[3]
Use Fourier transformation to solve formally

\[
-u''+u=f
\]

on \(\R\).
\end{exercise}

\begin{exercise}[3]
Transform the heat equation

\[
u_t=\kappa u_{xx}
\]

into frequency space and solve the resulting ordinary differential
equation.
\end{exercise}

\begin{exercise}[3]
Show that the Fourier transform preserves the \(L^2\) norm under
Plancherel's Theorem.
\end{exercise}

\begin{exercise}[4]
Prove Fej\'er's Theorem for continuous periodic functions.
\end{exercise}

\begin{exercise}[4]
Study the Gibbs phenomenon for the periodic square wave and determine the
behavior of its Fourier partial sums near a jump.
\end{exercise}

\begin{exercise}[4]
Prove the Riemann--Lebesgue Lemma.
\end{exercise}

\begin{exercise}[4]
Study a proof of Plancherel's Theorem beginning with Schwartz functions.
\end{exercise}

\begin{exercise}[4]
Prove a Fourier inversion theorem for Schwartz functions.
\end{exercise}

\begin{exercise}[4]
Show that the Fourier transform maps

\[
\mathcal{S}(\R)
\]

into itself.
\end{exercise}

\begin{exercise}[4]
Derive the classical uncertainty inequality using the Cauchy--Schwarz
inequality and integration by parts.
\end{exercise}

\begin{exercise}[4]
Prove the Poisson summation formula under suitable regularity and decay
assumptions.
\end{exercise}

\begin{exercise}[4]
Investigate how Fourier methods solve the heat equation on a bounded
interval using sine or cosine series.
\end{exercise}

\begin{exercise}[4]
Compare Fourier series, Fourier transforms, and discrete Fourier
transforms in terms of their domains and frequency sets.
\end{exercise}

%%%%%%%%%%%%%%%%%%%%%%%%%%%%%%%%%%%%%%%%%%%%%%%%%%%%%%%%%%%%
\subsection{Conceptual Summary}
\label{subsec:fourier-analysis-summary}
%%%%%%%%%%%%%%%%%%%%%%%%%%%%%%%%%%%%%%%%%%%%%%%%%%%%%%%%%%%%

Fourier analysis decomposes functions into oscillatory components.

For periodic functions,

\[
\boxed{
f(x)
\sim
\sum_{n=-\infty}^{\infty}
c_ne^{inx}.
}
\]

The coefficients are

\[
\boxed{
c_n
=
\frac{1}{2\pi}
\int_{-\pi}^{\pi}
f(x)e^{-inx}\,dx.
}
\]

In \(L^2\), these coefficients are Hilbert-space coordinates.

Parseval's identity gives

\[
\boxed{
\int_{-\pi}^{\pi}|f(x)|^2\,dx
=
2\pi
\sum_{n\in\Z}|c_n|^2.
}
\]

For functions on the real line, the Fourier transform is

\[
\boxed{
\widehat{f}(\xi)
=
\int_{\R}
f(x)e^{-2\pi ix\xi}\,dx.
}
\]

Important correspondences include

\[
\boxed{
f(x-a)
\longleftrightarrow
e^{-2\pi ia\xi}\widehat{f}(\xi),
}
\]

\[
\boxed{
e^{2\pi iax}f(x)
\longleftrightarrow
\widehat{f}(\xi-a),
}
\]

\[
\boxed{
f'(x)
\longleftrightarrow
2\pi i\xi\widehat{f}(\xi),
}
\]

and

\[
\boxed{
f*g
\longleftrightarrow
\widehat{f}\widehat{g}.
}
\]

Plancherel's Theorem gives

\[
\boxed{
\|\widehat{f}\|_2
=
\|f\|_2.
}
\]

The uncertainty principle expresses a fundamental tradeoff between
localization in the original and frequency variables.

Fourier analysis therefore converts many questions involving oscillation,
differentiation, convolution, and differential equations into simpler
frequency-domain problems.

%%%%%%%%%%%%%%%%%%%%%%%%%%%%%%%%%%%%%%%%%%%%%%%%%%%%%%%%%%%%
\subsection{Section Connections}
\label{subsec:fourier-analysis-connections}
%%%%%%%%%%%%%%%%%%%%%%%%%%%%%%%%%%%%%%%%%%%%%%%%%%%%%%%%%%%%

\chapterconnection{
Fourier Analysis is a central meeting point of Analysis, Complex
Analysis, Measure Theory, Lebesgue Integration, and Functional Analysis.
Fourier series are orthogonal expansions in Hilbert spaces, while the
Fourier transform becomes a unitary operator on \(L^2\) through
Plancherel's Theorem. Convolution and approximate identities connect
Fourier methods to approximation and regularization. Differentiation
becomes multiplication in frequency space, making Fourier methods
fundamental for Differential Equations and Mathematical Physics. The
Schwartz space and Fourier transform lead naturally to Distribution
Theory. Estimates for Fourier transforms, singular integrals, maximal
operators, and function spaces are developed further in Harmonic
Analysis. Fourier methods also appear throughout Probability, Number
Theory, Numerical Analysis, Signal Processing, and Partial Differential
Equations.
}

%%%%%%%%%%%%%%%%%%%%%%%%%%%%%%%%%%%%%%%%%%%%%%%%%%%%%%%%%%%%
% SECTION SOURCE TRACKING
%%%%%%%%%%%%%%%%%%%%%%%%%%%%%%%%%%%%%%%%%%%%%%%%%%%%%%%%%%%%

%------------------------------------------------------------
% 5.10 Fourier Analysis
%------------------------------------------------------------
%
% Source basis:
%   - Standard undergraduate and beginning graduate Fourier analysis.
%   - Standard Hilbert-space treatment of Fourier series.
%   - Standard Fourier-transform theory on R and L^2.
%
% Used for:
%   - periodic functions and harmonics;
%   - trigonometric orthogonality;
%   - real Fourier series;
%   - Fourier series on general intervals;
%   - even and odd expansions;
%   - half-range sine and cosine series;
%   - complex Fourier series;
%   - Fourier coefficients as Hilbert-space coordinates;
%   - Bessel's inequality and Parseval's identity;
%   - L^2 and pointwise convergence;
%   - Dirichlet kernels;
%   - Gibbs phenomenon;
%   - Fejer means and Fejer's Theorem;
%   - periodic convolution;
%   - approximate identities;
%   - Fourier transforms;
%   - Fourier inversion;
%   - translation, modulation, and scaling;
%   - differentiation and multiplication identities;
%   - convolution on R;
%   - Young's inequality;
%   - Riemann--Lebesgue Lemma;
%   - Plancherel's Theorem;
%   - Schwartz functions;
%   - Gaussian transforms;
%   - uncertainty principle;
%   - Poisson summation;
%   - Fourier methods for differential equations;
%   - heat and wave equations;
%   - filtering;
%   - discrete Fourier transform preview;
%   - probability and functional-analysis connections.
%
% Notes:
%   - Trigonometric functions and complex exponentials are developed
%     earlier in the text.
%   - Measure Theory and Lebesgue Integration provide the L^p
%     framework used throughout this section.
%   - Functional Analysis supplies the Hilbert-space interpretation
%     of Fourier expansions.
%   - Complex Analysis supplies contour methods and analytic
%     techniques related to Fourier integrals.
%   - Harmonic Analysis develops deeper Fourier estimates and
%     singular-integral methods.
%   - Distribution Theory later extends the Fourier transform to
%     generalized functions.
%
%%%%%%%%%%%%%%%%%%%%%%%%%%%%%%%%%%%%%%%%%%%%%%%%%%%%%%%%%%%%